\section{Evaluation Method}
\subsection{Fixtures and Measurements}
Use a Participants panel sample containing fictional names, including duplicate
display names. Target 20 visible entries only if they remain legible in one
screenshot; otherwise record the actual visible count. Maintain an independently
checked expected roster. Compare raw OCR with that roster before correction,
reporting exact-name matches, missing and extra entries, and the number of manual
corrections. No accuracy threshold is claimed before an initial OCR trial.

For the comparisons motivated by the literature, keep the expected roster fixed
and vary one input condition at a time, such as cropping or readable image scale.
Preserve duplicate entries when matching outputs to expected rows. Record engine
and model versions, language settings, preprocessing, and browser or service
configuration. Separate initial model-loading time from subsequent recognition
time when relevant. Report each tested condition and failure rather than
generalizing from the best sample to all participant screenshots.

After review, let $N$ denote the number of approved entries. Require exactly
$N$ distinct forms and $N$ correct name-to-link mappings. Check copied questions
against the template and submit fictional responses to verify that each response
appears only in its corresponding form. Test responder access using a
representative intended respondent rather than relying on the owner's editor view.

\subsection{Failure Scenarios and Reproducibility}
Test missing, invalid, empty, and oversized uploads; cancellation; duplicate
names; and partial service failures. Interrupt runs after spreadsheet creation,
form copying, state saving, and row updates. Resuming the same run must complete
remaining work without creating additional forms. Record cleanup outcomes for
temporary OCR resources in both successful and failed runs.

For each evaluation, record the source commit, account configuration, template
settings, fixture identifier, approved entry count, commands or UI steps, and
actual outputs. Measure total runtime and report limitations. If time savings
are evaluated, compare manual and automated workflows on the same fixture,
including human correction time; do not infer savings from successful execution.

